\subsection{Prolog basics}
\label{sec:prolog-basics}
\slidesrc{Prolog I: 1--37}

\subsubsection{Predicates and facts}
\slidesrc{Prolog I: 4--5}

In most programming languages the basic abstraction is the \emph{function} (or method): a
mapping from inputs to outputs. In Prolog the basic abstraction is the \textbf{predicate}
(or \textbf{relation}): it indicates \emph{whether or not some entities are related} in a
particular way. In logic programming we may \emph{test} whether entities are related, or
\emph{find} entities that are related to other entities.

\begin{nbexample}{the royal family program (slide 5)}
\begin{minted}{prolog}
parent(elizabeth, charles).
parent(diana, william).
parent(diana, harry).
parent(philip, charles).
parent(charles, william).
parent(charles, harry).
\end{minted}
\begin{itemize}[nosep]
  \item \code{parent} is a \textbf{predicate} (of arity 2, written \code{parent/2});
  \item each line is a separate \textbf{clause}, terminated by a full stop;
  \item each clause is a separate \emph{true statement}, understandable on its own:
        ``Elizabeth is a parent of Charles'';
  \item names of individuals (\code{elizabeth}) are \textbf{atoms}: constants, starting
        with a lower-case letter.
\end{itemize}
A Prolog program may consist of many files, each defining many predicates, each defined by
many clauses.
\end{nbexample}

\begin{center}
\begin{forest}
for tree={draw, rounded corners, font=\small, l sep=6mm, s sep=4mm, edge={->, thick}, fill=nbblue!8}
[{elizabeth $+$ philip}, fill=nbpurple!10
  [charles
    [william, fill=nbgreen!10] [harry, fill=nbgreen!10]]]
\end{forest}
\hspace{1.5cm}
\begin{forest}
for tree={draw, rounded corners, font=\small, l sep=6mm, s sep=4mm, edge={->, thick}, fill=nbblue!8}
[{charles $+$ diana}, fill=nbpurple!10
  [william, fill=nbgreen!10] [harry, fill=nbgreen!10]]
\end{forest}
\captionof{figure}{The family encoded by the \code{parent/2} facts (arrows go from parent to
child).}
\end{center}

\subsubsection{Making queries}
\slidesrc{Prolog I: 6--7}

Prolog systems provide an interactive environment to load and test code. The prompt is
\code{?-}. A source file (extension \code{.pl}) is loaded by writing its name in single
quotes inside square brackets, followed by a full stop: \code{?- ['royals.pl'].} Queries are
typed at the prompt and Prolog answers them:
\begin{minted}[linenos=false]{prolog}
?- parent(diana, harry).
true.
?- parent(elizabeth, diana).
false.
\end{minted}
(The slides' tuProlog prints \code{Yes} and \code{No}.)

\begin{nbwarning}[``false'' means ``not provable'', not ``false in reality'']
\code{parent(elizabeth, diana)} fails because it cannot be derived from the program. Prolog
works under the \textbf{closed-world assumption}: what is not known to be true is treated
as false. Ask \code{?- parent(elizabeth, anne).} and you also get \code{false}, although
Princess Anne is Elizabeth's daughter: she is simply not in the program.
\end{nbwarning}

\subsubsection{Variables and multiple directions}
\slidesrc{Prolog I: 8--11}

\begin{minted}[linenos=false]{prolog}
?- parent(elizabeth, X).
X = charles.
\end{minted}
Read the query as ``find \code{X} such that Elizabeth is a parent of \code{X}''.

\begin{nbdef}{Prolog variables (slide 8)}
Prolog variables begin with a \textbf{capital letter or an underscore}, followed by letters,
digits and underscores (\code{X}, \code{Kid}, \code{\_tmp}). They are very different from
variables in other programming languages: like in Haskell, \textbf{you cannot change the value
of a variable once it is set}. A variable stands for an \emph{unknown} object, as in algebra.
\end{nbdef}

Some queries have more than one answer. Prolog prints one answer and waits: type \code{;}
to ask for the next one, Return to stop.
\begin{minted}[linenos=false]{prolog}
?- parent(diana, X).
X = william ;
X = harry.
\end{minted}
(The slides show \code{X = harry ;} followed by \code{No}: tuProlog only discovers that
there are no more answers after you ask; SWI-Prolog often knows it already and ends with a
full stop.)

\begin{nbfact}{Multiple directions (slides 10--11)}
The same predicate can be queried ``the other way around'':
\begin{minted}[linenos=false]{prolog}
?- parent(X, william).
X = diana ;
X = charles.
\end{minted}
``Find \code{X} such that \code{X} is a parent of William.'' This works because the program
is \textbf{declarative}: it specifies information without specifying how it must be used.
For each argument we may pass a value or a variable; values \emph{constrain} the solutions,
variables do not. \code{?- parent(X, Y).} imposes no constraint and enumerates all six
pairs. Well-written logic programs exhibit this flexibility, although it is not always
practical (see \code{rev/2} in Section~\ref{sec:meta-programming}).
\end{nbfact}

\begin{nbintuition}[A predicate is a table, a query is a database question]
\notinslides
\code{parent/2} is like a two-column database table. \code{parent(diana, X)} is
``\code{SELECT child WHERE parent = 'diana'}''; \code{parent(X, william)} is
``\code{SELECT parent WHERE child = 'william'}''. In Java you would need two methods
(\code{getChildren}, \code{getParents}); in Prolog one relation answers both questions.
\end{nbintuition}

\subsubsection{Compound queries and constraints}
\slidesrc{Prolog I: 12--18}

A query can invoke several predicates separated by commas (``and''). If a variable appears
more than once, all occurrences must have the \emph{same value}, as in algebra.
\begin{minted}[linenos=false]{prolog}
?- parent(elizabeth, X), parent(X, Y).
X = charles, Y = william ;
X = charles, Y = harry.
\end{minted}

\begin{nbexample}{how Prolog solves a compound query (backtracking)}
\notinslides
\begin{enumerate}[nosep]
  \item First goal \code{parent(elizabeth, X)}: scan the facts top-down; the first match is
        \code{parent(elizabeth, charles)}, so \code{X = charles}. Remember this choice point.
  \item Second goal, now \code{parent(charles, Y)}: first match
        \code{parent(charles, william)}, \code{Y = william}. Print the answer.
  \item You type \code{;}: Prolog \textbf{backtracks} to the most recent choice point, the
        second goal, and continues scanning: \code{parent(charles, harry)},
        \code{Y = harry}. Print.
  \item \code{;} again: no more \code{parent(charles, \_)} facts; back to the first goal:
        no other \code{parent(elizabeth, \_)} fact. No more answers.
\end{enumerate}
\end{nbexample}

Let us add gender information (slide 13):
\begin{minted}{prolog}
female(elizabeth).   female(diana).
male(philip).        male(charles).      male(william).     male(harry).
\end{minted}
Predicates can take any number of arguments; the number is the \textbf{arity}. A predicate
with one argument defines a \emph{property} that certain entities may have.

\begin{nbexample}{Motherhood (slides 14--16)}
Give a Prolog query that determines the mother of William.
\begin{minted}[linenos=false]{prolog}
?- parent(X, william), female(X).
X = diana.
?- female(X), parent(X, william).     % equally good
X = diana.
\end{minted}
In the first version Prolog tries \code{X = diana} (fact 2), which is female: answer. (On
\code{;} it would try \code{X = charles}, who is not female: no more answers.) In the
second, it tries \code{X = elizabeth} first (not a parent of William), then
\code{X = diana}: answer.
\end{nbexample}

\begin{nbfact}{Constraints (slides 17--18)}
Think of each part of a compound query as a \textbf{constraint} on the solution: Prolog only
shows solutions that satisfy all of them. Some constraints have no solutions; in general
constraints can be placed in any order, but some orders are more efficient than others, and
some orders may not work as expected (more later).
\begin{minted}[linenos=false]{prolog}
?- parent(X,Y), parent(Y,Z), male(Z), female(X).
X = elizabeth, Y = charles, Z = william ;
X = elizabeth, Y = charles, Z = harry ;
false.
?- parent(X,Y), parent(Y,X).
false.
?- male(X), female(X).
false.
\end{minted}
\end{nbfact}

\subsubsection{Rules: clause bodies}
\slidesrc{Prolog I: 19--21}

\begin{nbdef}{Clause, head, body, fact, rule (slide 19)}
A clause may have two parts, a \textbf{head} and a \textbf{body}, separated by \code{:-}
(read ``\textbf{if}''). The head looks like a predicate with arguments (some may be
variables), the body looks like a query. Meaning: \emph{the head is true if the body is
true}. If the body is empty, \code{:-} is omitted: a clause with no body is a \textbf{unit
clause} or \textbf{fact}. A clause with a body is a \textbf{rule}.
\end{nbdef}

\begin{minted}{prolog}
grandparent(G, C) :-
    parent(G, P),
    parent(P, C).
\end{minted}
Read: ``\code{G} is a grandparent of \code{C} \emph{if} \code{G} is a parent of \code{P}
\emph{and} \code{P} is a parent of \code{C}''. Variables in a rule are universally
quantified: in FOL, $\forall G\,\forall C\,\forall P\,(\mathit{parent}(G,P) \land
\mathit{parent}(P,C) \limp \mathit{grandparent}(G,C))$.

\begin{minted}[linenos=false]{prolog}
?- grandparent(elizabeth, X).
X = william ;
X = harry.
?- grandparent(X, william).
X = elizabeth ;
X = philip ;
false.
\end{minted}

\begin{center}
\begin{forest}
for tree={draw, rounded corners, font=\small\ttfamily, l sep=5mm, s sep=8mm, edge={thick}, align=center}
[{?- grandparent(elizabeth, X)}, fill=nbblue!10
  [{?- parent(elizabeth, P), parent(P, X)}, fill=nbblue!10, edge label={node[midway, right, font=\scriptsize\rmfamily]{rule, G = elizabeth, C = X}}
    [{?- parent(charles, X)}, fill=nbblue!10, edge label={node[midway, right, font=\scriptsize\rmfamily]{P = charles}}
      [{success\\X = william}, fill=nbgreen!15]
      [{success\\X = harry}, fill=nbgreen!15]]]]
\end{forest}
\captionof{figure}{The search tree of \code{?- grandparent(elizabeth, X).}: each node is the
list of goals still to prove; each edge is the choice of a clause. Answers are the leaves
where nothing is left to prove. Section~\ref{sec:logical-foundations} calls this an
\emph{SLD tree}.}
\end{center}

\begin{nbexercise}{P1.1}{Family relationships (slide 21)}
\begin{enumerate}[nosep]
  \item Define a predicate \code{father(Dad, Kid)} that holds when \code{Dad} is the father
        of \code{Kid}.
  \item Define a predicate \code{greatgrandparent(Adult, Child)} that holds when
        \code{Adult} is a great-grandparent of \code{Child}.
\end{enumerate}
\end{nbexercise}
\begin{nbsolution}
\inputminted{prolog}{sections/prolog/prolog-basics/assets/family-solutions.pl}
\code{?- father(X, harry).} gives \code{X = charles} (the parents of Harry are Diana and
Charles; only Charles is male). With the six facts of the program,
\code{greatgrandparent/2} has \emph{no} solutions (the family tree has only three
generations): \code{?- greatgrandparent(A, C).} answers \code{false}. Adding, e.g.,
\code{parent(william, george).} makes \code{?- greatgrandparent(A, george).} answer
\code{A = elizabeth ; A = philip}. \textbf{Lesson}: reuse existing predicates; a
great-grandparent is a parent of a grandparent.
\end{nbsolution}

\subsubsection{Recursive predicates}
\slidesrc{Prolog I: 22--23}

\begin{minted}{prolog}
ancestor(A, A).
ancestor(A, D) :-
    parent(A, C),
    ancestor(C, D).
\end{minted}
The first clause says that everyone is their own ancestor (not the usual definition, but
common in computer science); the second that the parents of your ancestors are your
ancestors. The same code finds both ancestors and descendants:
\begin{minted}[linenos=false]{prolog}
?- ancestor(A, harry).
A = harry ; A = elizabeth ; A = diana ; A = philip ; A = charles ; false.
?- ancestor(elizabeth, D).
D = elizabeth ; D = charles ; D = william ; D = harry ; false.
\end{minted}

\begin{nbexample}{why the answers of \texttt{ancestor(A, harry)} come in that order}
\notinslides
\begin{enumerate}[nosep]
  \item Clause 1 unifies \code{ancestor(A, A)} with \code{ancestor(A, harry)}: \code{A = harry}.
  \item Clause 2: prove \code{parent(A, C), ancestor(C, harry)}, trying the parent facts in order:
  \begin{itemize}[nosep]
    \item \code{parent(elizabeth, charles)}: need \code{ancestor(charles, harry)}; clause 1
          fails (\code{charles} $\neq$ \code{harry}); clause 2: \code{parent(charles, william)}
          then \code{ancestor(william, harry)}: fails (William has no children);
          \code{parent(charles, harry)} then \code{ancestor(harry, harry)}: clause 1 succeeds.
          Answer \code{A = elizabeth}.
    \item \code{parent(diana, william)}: \code{ancestor(william, harry)} fails.
    \item \code{parent(diana, harry)}: \code{ancestor(harry, harry)} succeeds: \code{A = diana}.
    \item \code{parent(philip, charles)}: as for Elizabeth: \code{A = philip}.
    \item \code{parent(charles, william)}: fails; \code{parent(charles, harry)}: \code{A = charles}.
  \end{itemize}
\end{enumerate}
The order of answers follows the order of the facts, exactly as in a depth-first search.
\end{nbexample}

\begin{nbwarning}[Put the base case first and recurse on something smaller]
If the body were written \code{ancestor(C, D), parent(A, C)} (recursive call first), the
query \code{?- ancestor(A, harry).} would find the first answers and then call
\code{ancestor} with both arguments unknown, again and again, without ever consuming a
\code{parent} fact: an infinite loop (\emph{left recursion}). Same logic, different
behaviour: the order of goals matters for termination (Section~\ref{sec:logical-foundations}).
\end{nbwarning}

\subsubsection{Disjunction}
\slidesrc{Prolog I: 24}

Disjunction (``or'') in Prolog uses the \code{;} operator; it is equivalent to writing two
clauses:
\begin{center}
\begin{minipage}{0.42\linewidth}
\begin{minted}[linenos=false]{prolog}
person(X) :-
    ( good(X) ; bad(X) ).
\end{minted}
\end{minipage}
\quad$\equiv$\quad
\begin{minipage}{0.42\linewidth}
\begin{minted}[linenos=false]{prolog}
person(X) :- good(X).
person(X) :- bad(X).
\end{minted}
\end{minipage}
\end{center}

\subsubsection{Terms and variables}
\slidesrc{Prolog I: 25--28}

How can a variable be a data structure? A Prolog variable is \emph{just a term whose value
you don't know yet}. Don't \emph{assign} to variables, \emph{constrain} them. A variable can
be bound to any term, even another variable, but cannot change value once bound. It is
more like a variable in algebra than in a conventional programming language.

\begin{nbdef}{Compound term (slide 26)}
The \textbf{compound term} is Prolog's only data-structuring abstraction. It has a
\textbf{functor}, which is an atom, and one or more \textbf{arguments}, which can be any
terms. Syntax: functor first, then the arguments in parentheses, separated by commas. The
number of arguments is the \textbf{arity}. No need to declare it: just use it.
\end{nbdef}

Compound terms look just like predicate invocations (Whiteboard~\ref{wb:terms-atoms}: the
same syntax builds terms and atoms), but they are \emph{data}. Use them where you would use
(a pointer to) a \code{struct} in C; the functor indicates the kind of term:
\begin{center}
\code{position(20, 10)} \qquad \code{employee(1234, 'Jones', 'James', 100000)}
\end{center}
No allocation, no assignment to members: just write it down. Terms may contain variables,
whether or not their values are determined: \code{position(X, Y)},
\code{employee(Num, Surname, Firstname, Salary)}.

\begin{center}
\small
\begin{tabular}{@{}lll@{}}
\toprule
\textbf{Kind of term} & \textbf{Examples} & \textbf{Rule} \\
\midrule
atom & \code{elizabeth}, \code{[]}, \code{'Jones'}, \code{+} & lower-case start, or quoted \\
number & \code{42}, \code{-3}, \code{3.14} & \\
variable & \code{X}, \code{Kid}, \code{\_}, \code{\_tmp} & upper-case or \code{\_} start \\
compound & \code{position(20,10)}, \code{s(s(0))}, \code{[1,2]} & functor(args) \\
\bottomrule
\end{tabular}
\captionof{table}{The kinds of Prolog terms (added). Atoms and numbers are called
\emph{atomic}; a term with no variables is \emph{ground}.}
\end{center}

\subsubsection{Unification}
\slidesrc{Prolog I: 29--33}

In Prolog, \code{=} is used both for giving values to variables and for comparing terms:
\begin{minted}[linenos=false]{prolog}
?- 4 = 4.            true.
?- 4 = 5.            false.
?- X = 5.            X = 5.
?- 5 = V.            V = 5.
?- X = Y.            X = Y.
\end{minted}
(For the last query the slides show \code{X = \_G168, Y = \_G168}: both variables are bound
to the same fresh variable.) Variables appearing in different queries are unrelated.

\begin{nbdef}{Unification (slide 30)}
\code{=} performs \textbf{unification}: it tries to make two terms the \emph{same} by binding
the variables appearing in them. It is a powerful sort of pattern matching, and Prolog also
uses it for argument passing (calling \code{parent(diana, X)} unifies the call with the
head of each clause).
\begin{minted}[linenos=false]{prolog}
?- foo(A,2,C,D) = foo(1,X,Y,X).
A = 1, X = 2, D = 2, C = Y.
?- bar(X, X) = bar(1, 2).
false.
\end{minted}
\end{nbdef}

\begin{nbfact}{How to unify two terms (the algorithm, added)}
\notinslides
To unify $s$ and $t$:
\begin{enumerate}[nosep]
  \item if one of them is an unbound variable, bind it to the other (if both are the same
        variable, do nothing);
  \item if both are atomic (atoms or numbers), they unify iff they are identical;
  \item if both are compound, they unify iff they have the \textbf{same functor and arity}
        and their arguments unify \textbf{pairwise}, left to right, carrying the bindings
        along;
  \item otherwise they do not unify.
\end{enumerate}
The result is the \textbf{most general unifier} (mgu): the bindings that make the terms
equal while committing to as little as possible.
\end{nbfact}

\begin{nbexample}{\texttt{foo(A,2,C,D) = foo(1,X,Y,X)} step by step}
Same functor \code{foo}, same arity 4: unify the arguments pairwise.
\begin{center}
\small
\begin{tabular}{@{}clll@{}}
\toprule
\textbf{Pair} & \textbf{Left} & \textbf{Right} & \textbf{Result} \\
\midrule
1 & \code{A} & \code{1} & bind \code{A = 1} \\
2 & \code{2} & \code{X} & bind \code{X = 2} \\
3 & \code{C} & \code{Y} & bind \code{C = Y} (two variables, now the same unknown) \\
4 & \code{D} & \code{X} & \code{X} is already \code{2}: bind \code{D = 2} \\
\bottomrule
\end{tabular}
\end{center}
For \code{bar(X, X) = bar(1, 2)}: pair 1 binds \code{X = 1}; pair 2 must unify
\code{1} with \code{2}: different atomic terms, failure.
\end{nbexample}

\begin{nbexercise}{P1.2}{Unification (slides 31--33)}
Find bindings that make these pairs of terms the same:
(1) \code{f(X, 17) = f(Y, Y)}; \quad (2) \code{f(42, U, a) = f(V, f(V,W), W)}.
(\code{f} is used with arity 2 and 3: no problem, they are different functors.)
\end{nbexercise}
\begin{nbsolution}
(1) Pair 1: \code{X = Y}. Pair 2: \code{17 = Y}, so \code{Y = 17}, hence \code{X = 17}.
Answer: \code{X = 17, Y = 17}.

(2) Pair 1: \code{42 = V}, so \code{V = 42}. Pair 2: \code{U = f(V, W)} $=$ \code{f(42, W)}.
Pair 3: \code{a = W}, so \code{W = a}, and therefore \code{U = f(42, a)}. Answer:
\code{V = 42, W = a, U = f(42, a)}.
\end{nbsolution}

\begin{nbwarning}[Unification is not assignment, and not arithmetic]
\code{X = 1 + 2} binds \code{X} to the \emph{term} \code{+(1,2)}, not to \code{3}
(arithmetic needs \code{is}, Section~\ref{sec:arithmetic-nondeterminism}). And
\code{X = f(X)}: standard unification would need the \emph{occurs check} (a variable cannot
be bound to a term containing itself) and fail; SWI-Prolog skips the check for speed and
builds a cyclic term. Avoid such goals.
\end{nbwarning}

\subsubsection{Recursive structures: numbers as terms}
\slidesrc{Prolog I: 34--37}

Any argument of a compound term can be any term, including a term with the same functor:
this is how Prolog implements recursive structures (lists, trees, \ldots). A good strategy to
write predicates on recursive structures:
\begin{enumerate}[nosep]
  \item write a predicate that \emph{recognises} the recursive data structure;
  \item use it as a ``skeleton'' for other predicates, adding arguments and goals.
\end{enumerate}

For example, represent a natural number as either \code{0} or \code{s(X)} (``one more than
\code{X}'', the \emph{successor}), where \code{X} is a natural number: $3$ is
\code{s(s(s(0)))}. (This is \emph{not} how Prolog really does maths.)
\begin{minted}{prolog}
nat(0).                     % 0 is a natural number
nat(s(X)) :- nat(X).        % s(X) is a nat if X is

add(0, Y, Y).               % adding 0 to Y gives Y
add(s(X), Y, s(Z)) :-       % adding X+1 to Y gives Z+1
    add(X, Y, Z).           %   where Z is X + Y
\end{minted}
The structure of \code{add/3} follows that of \code{nat/1}. We refer to predicates as
\code{name/arity} because Prolog allows different predicates with the same name and
different arity.

\begin{nbexample}{\texttt{add(s(s(0)), s(s(s(0))), X)}: $2 + 3$}
\begin{align*}
&\code{add(s(s(0)), s(s(s(0))), X)} && \reason{clause 2: \code{X = s(Z1)}} \\
&\code{add(s(0), s(s(s(0))), Z1)} && \reason{clause 2: \code{Z1 = s(Z2)}} \\
&\code{add(0, s(s(s(0))), Z2)} && \reason{clause 1: \code{Z2 = s(s(s(0)))}}
\end{align*}
Back-substituting: \code{X = s(Z1) = s(s(Z2)) = s(s(s(s(s(0)))))}, i.e.\ $5$. Check:
$2 + 3 = 5$.
\end{nbexample}

\begin{nbfact}{Example queries (slide 36)}
\begin{minted}[linenos=false, fontsize=\footnotesize]{prolog}
?- nat(s(s(0))).                         true.
?- nat(X).                               X = 0 ; X = s(0) ; X = s(s(0)) ; ...
?- add(s(s(0)), X, Y).                   Y = s(s(X)).
?- add(s(s(0)), X, s(s(s(s(s(0)))))).    X = s(s(s(0))).        % 5 - 2
?- add(X, s(s(s(0))), s(s(s(s(s(0)))))). X = s(s(0)) ;  false.  % 5 - 3
?- add(X, Y, s(s(0))).                   X = 0, Y = s(s(0)) ;
                                         X = s(0), Y = s(0) ;
                                         X = s(s(0)), Y = 0 ; false.
\end{minted}
The same definition of addition also computes \textbf{subtraction} and enumerates all ways
of writing a number as a sum: multiple directions again. \code{nat(X)} generates the natural
numbers forever.
\end{nbfact}

\begin{nbexercise}{P1.3}{Multiplication (slide 37)}
Define \code{mult(X, Y, Z)}, true when \code{X} times \code{Y} equals \code{Z}, using
\code{add/3}.
\end{nbexercise}
\begin{nbsolution}
\inputminted{prolog}{sections/prolog/prolog-basics/assets/mult-solution.pl}
This follows the skeleton of \code{nat/1} on the first argument: $0 \cdot Y = 0$ and
$(X+1)\cdot Y = X\cdot Y + Y$. Trace of $2 \cdot 3$: \code{mult(s(s(0)), s(s(s(0))), Z)}
calls \code{mult(s(0), \ldots, W1)}, which calls \code{mult(0, \ldots, W0)}, giving
\code{W0 = 0}; then \code{add(0, s(s(s(0))), W1)} gives $W_1 = 3$; then
\code{add(s(s(s(0))), s(s(s(0))), Z)} gives $Z = 6$ = \code{s(s(s(s(s(s(0))))))}.
\end{nbsolution}
